\section{A bias-scale refinement with Gaussian-order sample exponents}
The preceding isotropic result permits trainable hidden biases at the
same scale as the other coordinates. Its sufficient pretraining exponent
depends on $p$. A simple variation of the prior removes this dependence.
This is a distinct anisotropic prior, not an improvement claimed for the
unchanged isotropic model.

First consider networks with hidden biases fixed to zero. H\"older gives
\[
 P:=\E_i|a_i|\|w_i\|_2
 \le r^{1/2-1/p}(V/2)^{2/p}
 \le(1+\epsilon)^{2/p}s/\rho
 \le(1+\epsilon)s/\rho.                              \tag{16}
\]
The path-norm form of the Gaussian transfer companion's deterministic
lemma therefore gives comparator error
$O_{p,k,s}(r^{2k+1}\sqrt{\epsilon+h/s})$, with the same norm and radius
orders as (6). Unlike the direct incoming-power projection bound, this
argument does not lose a $p$th root of the cost deficit.

The same conclusion holds with a positive continuous bias prior.
For $p>2$ put
\[
 \gamma=\min\{1,r\epsilon\}^{(p-2)/(2p)},
 \qquad
 V_\gamma=\E_i\left(|a_i|^p+|b_i/\gamma|^p
                                      +\sum_l|w_{i,l}|^p\right).
 \tag{17}
\]
At $p=2$ set $\gamma=1$. Give $a_i,w_{i,l}$ density (7) and give
$b_i/\gamma$ that same density, independently. Thus all physical
coordinates have continuous full-support priors, but the hidden bias
has scale $\gamma\tau$. The sharp minimum cost is still $V_*$:
the lower bound from the minimum-coordinate inequality is unchanged,
and the critical zero-bias witness attains it.

For $p>2$, weighted H\"older gives the incoming Euclidean bound
\[
 \|(b,w)\|_2\le
 \left(r+\gamma^{2p/(p-2)}\right)^{1/2-1/p}
       \left(|b/\gamma|^p+\sum_l|w_l|^p\right)^{1/p}.
\]
Consequently, if $V_\gamma\le(1+\epsilon)V_*$, then
\begin{align*}
 P&\le\left(r+\gamma^{2p/(p-2)}\right)^{1/2-1/p}
                                            (V_\gamma/2)^{2/p}\\
 &\le\frac{s}{\rho}(1+\epsilon)^{1/2-1/p+2/p}
 \le(1+\epsilon)\frac{s}{\rho}.                      \tag{18}
\end{align*}
Here $\gamma^{2p/(p-2)}/r\le\epsilon$, and
$1/2+1/p\le1$. At $p=2$ the same conclusion is the ordinary
quadratic arithmetic--geometric mean bound. The path-norm transfer
lemma again applies. The coefficient and radius bounds in (5) still
hold with $V_\gamma$ in place of $V$, since $0<\gamma\le1$.

\begin{theorem}[Polynomial transfer with a narrower continuous bias prior]
Fix $p\ge2$, $k$, and $s>0$. In the independent prior (17), use the
full-margin witness threshold (8), prior-scale bound (11), and IID
pretraining condition (12), with cost $V_\gamma$. Choose
\[
 \epsilon=\Theta_{p,k,s}(\delta^2r^{-4k-2})
\]
with a sufficiently small positive constant. Then an actual posterior
feature draw, followed by the same even-feature hinge fit, has error
$O(\delta)$ at the stated joint high probability. Sufficient orders are
\[
 N_0=O_{p,k,s}(r^{8k+7}\delta^{-4}),\quad
 m=O_{p,k,s}(r^{8k+7}\delta^{-4}\log(2/\zeta)),\quad
 n=O_{p,k,s}(r^{2k+3}\delta^{-2}\log(2/\zeta')).       \tag{19}
\]
These exponents no longer depend on $p$. Both physical prior scales can
be chosen inverse polynomially small. The theorem also holds with biases
fixed to zero, using $d=r+1$ scalar parameters per neuron in the evidence
calculation.
\end{theorem}
\begin{proof}
Use the same full-margin witness and boxes, with the bias-box coordinate
interpreted as $b/\gamma$. The witness has zero bias, its cost and mean
output are unchanged, and a bias-box perturbation changes the physical
bias by at most the original box radius. Thus all cost, strict-margin,
and evidence estimates (8)--(11) remain valid. The prior gamma moment
uses the $d$ standardized independent coordinates and is unchanged.
The common variation-ball sample event (12) also remains valid: the
physical variation norm is bounded either by its original expression or
by the sharper (18). The path-norm lemma, (18), and the original
readout/radius moment bounds give the asserted deterministic error and
downstream sample count. Substituting the new $\epsilon$ into
$N_0,m=O_{p,s}(r^3\epsilon^{-2})$ gives (19).
The previously proved lower bound on the admissible $\tau^p$ is inverse
polynomial, and (17) makes $\gamma$ inverse polynomial as well.
\end{proof}

This refinement concerns statistical sufficiency under an explicitly
specified prior. It does not prove that changing the bias scale makes
posterior sampling efficient, or that the isotropic $p>2$ exponents in
(13) are necessary. Sharper stability for the latter prior remains open.
